\papertitle{基于方向等效坡度与有限曲率优化的多波束测线规划}

\begin{cnabstract}
多波束测深的核心矛盾在于：测线过疏会造成漏测，测线过密又会增加重复覆盖和航行成本。本文围绕理想斜坡与真实海底地形下的覆盖计算和测线设计，依次建立二维扇形声束几何模型、任意航向等效坡度模型、最大允许间距递推模型和地形自适应有限曲率路线模型。

针对问题一，在垂直测线的斜坡剖面内分别求得深水侧和浅水侧覆盖半宽，并以当前条带宽度为基准定义有向重叠率。计算表明，水深由 \SI{90.95}{m} 降至 \SI{49.05}{m} 时，覆盖宽度由 \SI{315.81}{m} 降至 \SI{170.33}{m}；固定 \SI{200}{m} 间距下，浅水侧 \SI{600}{m} 与 \SI{800}{m} 位置的重叠率分别为 $-1.53\%$ 和 $-12.36\%$，已出现漏测。针对问题二，将坡度分解到航向与横航向，利用沿航向深度变化和横航向等效坡度构造三维覆盖宽度模型，得到 $8$ 个方向、$8$ 个距离下的完整矩阵；结果满足 $45^\circ$ 与 $315^\circ$、$135^\circ$ 与 $225^\circ$ 的对称关系，且 $90^\circ$、$270^\circ$ 方向的水深沿程不变。

针对问题三，将候选限制为平行等深线的南北向直线族，以西边界相切确定首线，以相邻重叠率取下限 $10\%$ 确定后续最大间距，并以首次覆盖东边界为终止条件。求得 $34$ 条测线，总长度为 \SI{125936}{m}；在同一规则下，$33$ 条测线仍留下 \SI{25.756620}{m} 的东侧缺口，因此 $34$ 条为该路线族内的最少条数。针对问题四，读取 $251\times201$ 深度网格，采用双线性插值与局部梯度计算任意位置的横航向等效坡度，构造独立空间覆盖评估器，并比较南北保守直线、东西保守直线、地形自适应折线和有限曲率平滑曲线。最终推荐 $61$ 条地形自适应有限曲率测线，几何总长度为 \SI{576726.636}{m}，名义网格漏测率为 $0$，重叠率超过 $20\%$ 的累计测线长度为 \SI{339556.407}{m}，最小曲率半径为 \SI{225.983}{m}。在 $70\times70$ 敏感性网格上出现 $0.040816\%$ 的微小漏测，故本文不作连续全覆盖或全局最优声明，并保留总长度 \SI{620420}{m} 的保守直线方案作为工程回退。

所建模型将解析推导、独立复算、结构对称性、边界检查和网格敏感性结合起来，能够直接回答四个问题，同时明确结论适用范围和剩余风险。
\end{cnabstract}

\keywords{多波束测深；方向等效坡度；最大间距递推；有限曲率路线；敏感性分析}

\clearpage
